\chapter{FUNCTIONAL ANALYSIS AND MODELING}

Functional analysis and modeling represent a critical phase in the system development life cycle, serving as the formal bridge that connects raw stakeholder requirements with concrete architectural designs, database schemas, and process flows. In the context of the proposed AI-Driven Lab Exam Monitoring and Management System, this phase is exceptionally crucial because of the complex, real-time coordination required between multiple system layers, including low-level Windows API hooks, secure background telemetry streams, real-time teacher dashboards, automated grading engines, and secure PDF reporting systems. By decomposing these requirements into visual and tabular models, the project team can systematically analyze system behavior, clarify boundaries, ensure security compliance, and verify that the proposed software prototype addresses both academic governance and technical safety.

This chapter is organized into two main sections: Use Case Modeling and Functional Modeling. The Use Case Modeling section focuses on the user-centric requirements of the system, defining the interactions between external actors and system functionalities. This section includes detailed user stories, a UML use case diagram, and individual actor use case specifications. The Functional Modeling section outlines the structural and behavioral aspects of the system, detailing the Entity Relationship Diagram (ERD) for data storage and the Data Flow Diagrams (DFDs) at various levels of abstraction to map the movement and transformation of telemetry and evaluation data.

\section{Use Case Modeling}

Use case modeling is a widely adopted Software Engineering technique for documenting functional requirements from the perspective of the system's end users \cite{uml_systems_analysis}. It identifies the external entities that interact with the system, known as actors, and the specific goals that these actors want to achieve, known as use cases. By focusing on what the system does rather than how it does it, use case modeling provides a clear, high-level map of the system's scope, boundaries, and primary interactions. This user-centric representation helps eliminate ambiguity during development and serves as a baseline for validating and testing system behaviors.

For the AI-Driven Lab Exam Monitoring and Management System, the use case model captures the interactions of three primary human actors: the Student, the Teacher, and the System Administrator, along with one automated external entity, the AI evaluation service. The student interacts with the workstation desktop client to run pre-exam diagnostics, receive authorization, complete tests under lockdown, and submit source code or theory files. The teacher utilizes the web dashboard to configure exams, define monitoring rules, observe live lab feeds, handle system overrides, verify AI grading metrics, and download reports. The administrator manages system configurations, updates user accounts and role assignments, and audits system logs. These interactions are illustrated in the UML Use Case Diagram shown in Figure~\ref{fig:usecase_diagram}.

The diagram presented in Figure~\ref{fig:usecase_diagram} maps the operational boundaries of the system. It demonstrates how student actors are confined to client-side activities such as hardware checks, identity validation, locked exam execution, and response submission. In contrast, teacher actors are granted access to management and monitoring use cases, such as exam creation, allowed applications list configuration, live telemetry observation, and review of automated evaluations. The administrator actor spans both client and server domains by managing system settings, managing user accounts, and conducting security audits. This visual separation of responsibilities ensures that authorization constraints are defined at the modeling level and are consistently enforced throughout implementation.

\begin{figure}[H]
\centering
\includegraphics[width=0.97\linewidth]{Figures/use_case_diagram.png}
\caption{Use Case Diagram of the Proposed System}
\label{fig:usecase_diagram}
\end{figure}

\subsection{User Stories}

User stories are a foundational practice in Agile software development, providing a lightweight, natural-language description of system functionality from the perspective of an end-user \cite{user_stories_applied}. Each story outlines the role of the user, the specific action they need to perform, and the underlying value or benefit they expect to receive. This structure ensures that development efforts remain focused on delivering actual value to stakeholders and makes it easy to discuss system features with non-technical users. In the development of the proposed system, these user stories were continuously refined and organized into sprints, serving as the basis for tracking progress and defining acceptance criteria.

The user stories for the AI-Driven Lab Exam Monitoring and Management System are organized by actor roles, covering the entire lifecycle of an examination from setup to monitoring, submission, evaluation, and reporting:

\subsubsection{Student User Stories}
\begin{enumerate}
    \item \textbf{AuthenticateUserSession:} As a student, I want to log in using my credentials, so that I can establish a secure, authenticated session.
    \item \textbf{VerifyWorkstationHardware:} As a student, I want the system to bind my workstation HWID, so that I can lock my exam session to this physical machine.
    \item \textbf{ExecutePreExamDiagnostics:} As a student, I want to run pre-exam system verification checks, so that I can resolve webcam, mic, or process compliance issues.
    \item \textbf{DisplayStudentDashboard:} As a student, I want to view my scheduled exams and read invigilator instructions, so that I can prepare to launch my assessment workspace.
    \item \textbf{LockdownStudentWorkstation:} As a student, I want the client to restrict task-switching, hotkeys, and unauthorized background applications, so that I can focus on the exam while demonstrating compliance.
    \item \textbf{ConductExamSession:} As a student, I want to start my exam, view questions, write answers, and test-compile code, so that I can complete the assessment inside the secure environment.
    \item \textbf{StreamActivityTelemetry:} As a student, I want the system to stream real-time screen captures and focus events, so that the teacher has continuous evidence of my exam integrity.
    \item \textbf{SubmitExamAnswers:} As a student, I want my answers to auto-save and to manually submit them, so that my work is secure and officially registered.
\end{enumerate}

\subsubsection{Teacher User Stories}
\begin{enumerate}
    \item \textbf{CreateExamTemplate:} As a teacher, I want to configure exam templates, questions, and test cases, so that the evaluation aligns with course standards.
    \item \textbf{ConfigureApplicationWhitelist:} As a teacher, I want to define whitelisted desktop software, so that students can access required programming or typing tools.
    \item \textbf{MapSeatingEligibility:} As a teacher, I want to map student registrations to laboratory seating and IP ranges, so that only scheduled candidates can access the exam from approved terminals.
    \item \textbf{MonitorLiveClassroom:} As a teacher, I want to monitor focus timelines and active processes, so that I can supervise the lab room from my dashboard.
    \item \textbf{EvaluateExamSubmissions:} As a teacher, I want the system to automatically compile code in Docker, semantic-grade theory answers, check plagiarism, and export PDF reports, so that I can grade submissions quickly.
\end{enumerate}

\subsubsection{Administrator User Stories}
\begin{enumerate}
    \item \textbf{ManageUserAccounts:} As an administrator, I want to manage user roles, system accounts, and configurations, so that system access complies with institutional policies.
\end{enumerate}

These user stories cover all 12 modules of the AI-Driven Lab Exam Monitoring and Management System. They establish the functional capabilities required of the software prototype, ensuring that the system is not merely an automated lockdown browser but a comprehensive exam administration and grading framework. By grouping these stories into development sprints, the project team can construct, test, and deliver features incrementally, validating that the student client and teacher web panel remain coordinated at every stage.

\subsection{Individual Actor Use Cases}

To refine the system's functional requirements, the user stories are expanded into formal individual actor use case specifications \cite{use_case_definition}. A use case specification provides a detailed, step-by-step description of how an actor interacts with the system to achieve a specific goal. It details the preconditions, triggers, main flow of events, postconditions, and exception flows, ensuring that developers, testers, and stakeholders share a common understanding of system behavior under both normal and anomalous conditions.

In order to maintain design consistency and align with institutional formatting requirements, every use case table in this document follows the exact double-hline layout (\texttt{\textbackslash hline\textbackslash hline}) and array stretch of 1.2. Below are the use case specifications for the fourteen functional areas of the SmartExam system:

In Table 4.1, the use case `ManageUserAccounts` allows administrators to configure user profiles, assign roles, and toggle account activity status. The preconditions require that the administrator is authenticated and the database is accessible. The trigger is the administrator selecting user account management on the dashboard. The postconditions verify that the user profile details are updated and recorded in the database, with exception flows handling duplicate emails or system locks.

\begin{table}[H]
\centering
\caption{Use Case Table of User Account Management (ManageUserAccounts)}
\label{tab:uct_manage_user_accounts}
\renewcommand{\arraystretch}{1.2}
\begin{tabular}{lp{10.6cm}}
\hline\hline
\textbf{Field} & \textbf{Details} \\
\hline\hline
\textbf{UCT ID} & UCT-01 \\
\textbf{Case Name} & ManageUserAccounts \\
\textbf{Overview} & Allows administrators to create, update, and manage user accounts and system permissions. \\
\textbf{Actors} & Administrator, System Database \\
\textbf{Pre-Conditions} & Administrator is authenticated and database connection is active. \\
\textbf{Triggers} & Administrator navigates to the User Management section. \\
\textbf{Steps} & \begin{minipage}[t]{10.6cm}
1. Admin enters account details (name, email, role) and submits.\\
2. System validates email format and uniqueness.\\
3. System writes user profile and default settings to the database.\\
4. System returns success message and updates the active user grid.
\end{minipage} \\
\textbf{Post-Conditions} & The user account is saved, and active roles are assigned. \\
\textbf{Exception Flow} & \begin{minipage}[t]{10.6cm}
1. Email Conflict: System flags duplicate email and rejects registration.\\
2. Validation Fail: System details invalid fields and prompts for corrections.
\end{minipage} \\
\hline\hline
\end{tabular}
\end{table}

In Table 4.2, the use case `AuthenticateUserSession` handles the user login process using credentials to verify identity. The preconditions require that the user is registered and the system is online. The trigger occurs when a user enters credentials on the login screen. The postconditions verify that the user has received a valid JWT token, and the exception flows prompt the user to re-enter details if authentication fails.

\begin{table}[H]
\centering
\caption{Use Case Table of User Session Authentication (AuthenticateUserSession)}
\label{tab:uct_authenticate_user_session}
\renewcommand{\arraystretch}{1.2}
\begin{tabular}{lp{10.6cm}}
\hline\hline
\textbf{Field} & \textbf{Details} \\
\hline\hline
\textbf{UCT ID} & UCT-02 \\
\textbf{Case Name} & AuthenticateUserSession \\
\textbf{Overview} & Authenticates student, teacher, or admin credentials and returns a secure session token. \\
\textbf{Actors} & Student, Teacher, Admin, System \\
\textbf{Pre-Conditions} & User is registered, client is running, and API server is online. \\
\textbf{Triggers} & User submits email and password on the login screen. \\
\textbf{Steps} & \begin{minipage}[t]{10.6cm}
1. User enters registered email and password.\\
2. Client transmits credentials securely to the server.\\
3. System verifies credentials against the database.\\
4. System generates a secure JWT token and returns it to the client.
\end{minipage} \\
\textbf{Post-Conditions} & User is authenticated and a session token is issued. \\
\textbf{Exception Flow} & \begin{minipage}[t]{10.6cm}
1. Invalid Credentials: System displays error and prompts retry.\\
2. Database Offline: System shows connection failure alert.
\end{minipage} \\
\hline\hline
\end{tabular}
\end{table}

In Table 4.3, the use case `VerifyWorkstationHardware` binds the workstation motherboard and CPU hashes to the student account. The preconditions require that the student has logged in and the WPF client is running. The trigger occurs automatically upon login. The postconditions ensure that the HWID is bound and logged, while exception flows lock the client if a device mismatch occurs.

\begin{table}[H]
\centering
\caption{Use Case Table of Workstation Hardware Verification (VerifyWorkstationHardware)}
\label{tab:uct_verify_workstation_hardware}
\renewcommand{\arraystretch}{1.2}
\begin{tabular}{lp{10.6cm}}
\hline\hline
\textbf{Field} & \textbf{Details} \\
\hline\hline
\textbf{UCT ID} & UCT-03 \\
\textbf{Case Name} & VerifyWorkstationHardware \\
\textbf{Overview} & Retreives and validates the workstation Hardware ID (HWID) hash to prevent proxy exam attempts. \\
\textbf{Actors} & Student, System \\
\textbf{Pre-Conditions} & User session is authenticated, and WPF client is executing. \\
\textbf{Triggers} & Authentication completes, initiating device check. \\
\textbf{Steps} & \begin{minipage}[t]{10.6cm}
1. WPF client queries local system APIs to fetch CPU, motherboard, and MAC address details.\\
2. Client hashes details to construct a unique Hardware ID (HWID) payload.\\
3. Client transmits HWID payload to the REST API server.\\
4. System compares HWID against registered bindings in the database.
\end{minipage} \\
\textbf{Post-Conditions} & Workstation HWID is validated or registered for binding. \\
\textbf{Exception Flow} & \begin{minipage}[t]{10.6cm}
1. HWID Mismatch: System blocks exam access and alerts invigilator.\\
2. System API Failure: WPF client details hardware querying error.
\end{minipage} \\
\hline\hline
\end{tabular}
\end{table}

In Table 4.4, the use case `CreateExamTemplate` allows teachers to add question items, configure duration settings, and establish grading parameters. The preconditions require that the teacher is authenticated and has course management privileges. The trigger is the teacher clicking the exam creation option. The postconditions ensure that the exam template is published and stored in the database.

\begin{table}[H]
\centering
\caption{Use Case Table of Exam Template Configuration (CreateExamTemplate)}
\label{tab:uct_create_exam_template}
\renewcommand{\arraystretch}{1.2}
\begin{tabular}{lp{10.6cm}}
\hline\hline
\textbf{Field} & \textbf{Details} \\
\hline\hline
\textbf{UCT ID} & UCT-04 \\
\textbf{Case Name} & CreateExamTemplate \\
\textbf{Overview} & Enables teachers to define exam parameters, questions, marks, and test configurations. \\
\textbf{Actors} & Teacher, System \\
\textbf{Pre-Conditions} & Teacher is authenticated with appropriate course access permissions. \\
\textbf{Triggers} & Teacher selects the "Create Exam" option on the dashboard. \\
\textbf{Steps} & \begin{minipage}[t]{10.6cm}
1. Teacher inputs exam title, instructions, and duration.\\
2. Teacher adds individual coding or theory question details.\\
3. Teacher uploads test cases for automated compilation checks.\\
4. Teacher publishes the exam configuration.
\end{minipage} \\
\textbf{Post-Conditions} & The exam template is created and stored in the database. \\
\textbf{Exception Flow} & \begin{minipage}[t]{10.6cm}
1. Empty Fields: System displays errors and requests missing details.\\
2. Invalid Test Cases: System flags compile-syntax errors in code inputs.
\end{minipage} \\
\hline\hline
\end{tabular}
\end{table}

In Table 4.5, the use case `ConfigureApplicationWhitelist` allows teachers to define which local computer programs are permitted during the exam. The preconditions require that the exam template is active. The trigger occurs when whitelisted application editing is initiated. The postconditions update whitelists in the database, with exception flows prompting correction for empty process formats.

\begin{table}[H]
\centering
\caption{Use Case Table of Whitelisted Applications Configuration (ConfigureApplicationWhitelist)}
\label{tab:uct_configure_application_whitelist}
\renewcommand{\arraystretch}{1.2}
\begin{tabular}{lp{10.6cm}}
\hline\hline
\textbf{Field} & \textbf{Details} \\
\hline\hline
\textbf{UCT ID} & UCT-05 \\
\textbf{Case Name} & ConfigureApplicationWhitelist \\
\textbf{Overview} & Defines whitelisted desktop programs allowed to execute during the exam session. \\
\textbf{Actors} & Teacher, System \\
\textbf{Pre-Conditions} & The targeted exam template exists and is customizable. \\
\textbf{Triggers} & Teacher selects the whitelisting menu on the configuration panel. \\
\textbf{Steps} & \begin{minipage}[t]{10.6cm}
1. Teacher adds application filenames (e.g., code.exe, python.exe) to the list.\\
2. Teacher specifies permitted paths or arguments if applicable.\\
3. Teacher saves the whitelist configuration.\\
4. System updates the exam rules table.
\end{minipage} \\
\textbf{Post-Conditions} & Whitelist settings are bound to the exam template. \\
\textbf{Exception Flow} & \begin{minipage}[t]{10.6cm}
1. Empty Whitelist: System prompts confirmation for zero whitelisted tools.\\
2. Format Mismatch: System flags invalid application extension strings.
\end{minipage} \\
\hline\hline
\end{tabular}
\end{table}

In Table 4.6, the use case `MapSeatingEligibility` maps registered students to physical lab terminals and network IP filters. The preconditions require that the student list is imported and the exam is scheduled. The trigger is the teacher configuring lab seating mapping. The postconditions ensure that seating assignments are active in the database.

\begin{table}[H]
\centering
\caption{Use Case Table of Seating and Eligibility Mapping (MapSeatingEligibility)}
\label{tab:uct_map_seating_eligibility}
\renewcommand{\arraystretch}{1.2}
\begin{tabular}{lp{10.6cm}}
\hline\hline
\textbf{Field} & \textbf{Details} \\
\hline\hline
\textbf{UCT ID} & UCT-06 \\
\textbf{Case Name} & MapSeatingEligibility \\
\textbf{Overview} & Maps eligible candidate logins to lab terminals, seat IDs, and permitted network ranges. \\
\textbf{Actors} & Teacher, System \\
\textbf{Pre-Conditions} & Student student profiles and exam records are available. \\
\textbf{Triggers} & Teacher accesses the seating mapping grid on the console. \\
\textbf{Steps} & \begin{minipage}[t]{10.6cm}
1. Teacher assigns student emails to specific seat coordinates.\\
2. Teacher enters allowed IP address filters for the exam room.\\
3. Teacher saves the mapped layout to lock terminal constraints.\\
4. System writes seating and network policies to the database.
\end{minipage} \\
\textbf{Post-Conditions} & Student login bounds are mapped to laboratory seats. \\
\textbf{Exception Flow} & \begin{minipage}[t]{10.6cm}
1. Overlapping Mappings: System rejects duplicate seat assignments.\\
2. Invalid IP Range: System details formatting errors and requests correction.
\end{minipage} \\
\hline\hline
\end{tabular}
\end{table}

In Table 4.7, the use case `ExecutePreExamDiagnostics` tests student workstation hardware, processes, and connection bandwidth. The preconditions require that the student is logged in and the client is ready. The trigger occurs when the student starts diagnostics. The postconditions generate passing diagnostic logs, while exception flows block the exam launch button.

\begin{table}[H]
\centering
\caption{Use Case Table of Pre-Exam Diagnostics (ExecutePreExamDiagnostics)}
\label{tab:uct_execute_pre_exam_diagnostics}
\renewcommand{\arraystretch}{1.2}
\begin{tabular}{lp{10.6cm}}
\hline\hline
\textbf{Field} & \textbf{Details} \\
\hline\hline
\textbf{UCT ID} & UCT-07 \\
\textbf{Case Name} & ExecutePreExamDiagnostics \\
\textbf{Overview} & Tests student hardware, webcam access, microphone levels, and checks for running blacklisted processes. \\
\textbf{Actors} & Student, System \\
\textbf{Pre-Conditions} & Student is logged in and the scheduled exam session is active. \\
\textbf{Triggers} & Student initiates the pre-exam system verification checks. \\
\textbf{Steps} & \begin{minipage}[t]{10.6cm}
1. WPF client checks webcam and microphone status.\\
2. Client scans the active process list for blacklisted software.\\
3. Client performs network latency and connection bandwidth checks.\\
4. Client prompts student to take a face baseline picture via webcam.\\
5. Client uploads diagnostic reports to the REST API.
\end{minipage} \\
\textbf{Post-Conditions} & Workstation diagnostic logs are generated and recorded. \\
\textbf{Exception Flow} & \begin{minipage}[t]{10.6cm}
1. Hardware Missing: System flags missing camera and disables exam entry.\\
2. Blacklisted Process: System lists running apps (e.g. Discord) and demands closure.
\end{minipage} \\
\hline\hline
\end{tabular}
\end{table}

In Table 4.8, the use case `DisplayStudentDashboard` displays student examination details, schedules, and start options. The preconditions require that the student has completed diagnostics. The trigger is the client redirecting to the dashboard. The postconditions load exam instructions, and exception flows display errors for inactive exam states.

\begin{table}[H]
\centering
\caption{Use Case Table of Student Dashboard Display (DisplayStudentDashboard)}
\label{tab:uct_display_student_dashboard}
\renewcommand{\arraystretch}{1.2}
\begin{tabular}{lp{10.6cm}}
\hline\hline
\textbf{Field} & \textbf{Details} \\
\hline\hline
\textbf{UCT ID} & UCT-08 \\
\textbf{Case Name} & DisplayStudentDashboard \\
\textbf{Overview} & Loads student dashboard showing active exam listings, instructions, and workspace launcher. \\
\textbf{Actors} & Student, System \\
\textbf{Pre-Conditions} & Pre-exam diagnostics checks have passed successfully. \\
\textbf{Triggers} & System redirects student after verification checks. \\
\textbf{Steps} & \begin{minipage}[t]{10.6cm}
1. Dashboard queries REST API for scheduled exam records.\\
2. REST API queries the database and returns exam details.\\
3. Dashboard displays instructions, duration, and question metadata.\\
4. Dashboard enables the "Launch Exam Workspace" button.
\end{minipage} \\
\textbf{Post-Conditions} & Exam instructions and dashboard layout are displayed. \\
\textbf{Exception Flow} & \begin{minipage}[t]{10.6cm}
1. Exam Not Started: Dashboard shows timer countdown to scheduled start.\\
2. Expired Session: Client prompts user to re-authenticate.
\end{minipage} \\
\hline\hline
\end{tabular}
\end{table}

In Table 4.9, the use case `LockdownStudentWorkstation` locks student workstations to fullscreen, disables hotkeys, and intercepts task-switching. The preconditions require that the student clicks launch exam. The trigger is the exam starting. The postconditions lock down the desktop environment, and exception flows sound alarms for hook failures.

\begin{table}[H]
\centering
\caption{Use Case Table of Student Workstation Lockdown (LockdownStudentWorkstation)}
\label{tab:uct_lockdown_student_workstation}
\renewcommand{\arraystretch}{1.2}
\begin{tabular}{lp{10.6cm}}
\hline\hline
\textbf{Field} & \textbf{Details} \\
\hline\hline
\textbf{UCT ID} & UCT-09 \\
\textbf{Case Name} & LockdownStudentWorkstation \\
\textbf{Overview} & Restricts desktop access, blocks Alt+Tab/Task Manager hooks, and restricts user inputs. \\
\textbf{Actors} & Student, System \\
\textbf{Pre-Conditions} & Scheduled exam session is started by the student. \\
\textbf{Triggers} & Student launches the active exam workspace environment. \\
\textbf{Steps} & \begin{minipage}[t]{10.6cm}
1. WPF client shifts to fullscreen mode covering all displays.\\
2. Client activates low-level Windows keyboard hooks to intercept hotkeys.\\
3. Client disables secondary monitor outputs if detected.\\
4. Client blocks Task Manager execution.
\end{minipage} \\
\textbf{Post-Conditions} & Workstation environment is locked down and restricted. \\
\textbf{Exception Flow} & \begin{minipage}[t]{10.6cm}
1. Hook Blocked: Client logs violation, sounds alarm, and terminates exam session.\\
2. Multiple Screens: Client restricts window view and prompts user to disconnect monitors.
\end{minipage} \\
\hline\hline
\end{tabular}
\end{table}

In Table 4.10, the use case `ConductExamSession` renders active exam questions, processes user inputs, and supports code compilations. The preconditions require that the workstation is locked. The trigger occurs when the workspace loads questions. The postconditions record answer entries in the cache, and exception flows cache answers locally if connections fail.

\begin{table}[H]
\centering
\caption{Use Case Table of Active Exam Session (ConductExamSession)}
\label{tab:uct_conduct_exam_session}
\renewcommand{\arraystretch}{1.2}
\begin{tabular}{lp{10.6cm}}
\hline\hline
\textbf{Field} & \textbf{Details} \\
\hline\hline
\textbf{UCT ID} & UCT-10 \\
\textbf{Case Name} & ConductExamSession \\
\textbf{Overview} & Renders exam questions, receives code or text answers, and runs student code tests. \\
\textbf{Actors} & Student, System \\
\textbf{Pre-Conditions} & Workstation lockdown is fully enforced and server connection is active. \\
\textbf{Triggers} & Workspace interface loads questions. \\
\textbf{Steps} & \begin{minipage}[t]{10.6cm}
1. Workspace displays code editors and text input zones.\\
2. Student types answers or completes programming exercises.\\
3. Student runs mock compilation to test code against public test cases.\\
4. System compiles and returns output results inside the client.
\end{minipage} \\
\textbf{Post-Conditions} & Student responses are structured and ready for submission. \\
\textbf{Exception Flow} & \begin{minipage}[t]{10.6cm}
1. Compilation Crash: Editor returns traceback details and errors.\\
2. Server Timeout: Client caches student answer inputs locally.
\end{minipage} \\
\hline\hline
\end{tabular}
\end{table}

In Table 4.11, the use case `StreamActivityTelemetry` collects process logs and screenshots and sends them to WebSocket hubs. The preconditions require that the exam session is active. The trigger is the workspace launching background telemetry threads. The postconditions store telemetry events in the database, with exception flows caching events.

\begin{table}[H]
\centering
\caption{Use Case Table of Real-Time Telemetry Streaming (StreamActivityTelemetry)}
\label{tab:uct_stream_activity_telemetry}
\renewcommand{\arraystretch}{1.2}
\begin{tabular}{lp{10.6cm}}
\hline\hline
\textbf{Field} & \textbf{Details} \\
\hline\hline
\textbf{UCT ID} & UCT-11 \\
\textbf{Case Name} & StreamActivityTelemetry \\
\textbf{Overview} & Streams screenshots, active focus events, and running processes over SignalR WebSockets. \\
\textbf{Actors} & Student, System WebSocket Server \\
\textbf{Pre-Conditions} & Student exam workspace is running. \\
\textbf{Triggers} & Background logging threads initiate telemetry collection. \\
\textbf{Steps} & \begin{minipage}[t]{10.6cm}
1. WPF client captures workstation screenshot at randomized intervals.\\
2. Client queries running processes and window focus handles.\\
3. Client bundles screenshot and process logs into JSON.\\
4. Client streams packet to the SignalR WebSocket hub.
\end{minipage} \\
\textbf{Post-Conditions} & Telemetry packets are written to the database logs. \\
\textbf{Exception Flow} & \begin{minipage}[t]{10.6cm}
1. Connection Lost: Client stores telemetry in encrypted local cache.\\
2. Camera Evasion: System flags missing face stream and fires warnings.
\end{minipage} \\
\hline\hline
\end{tabular}
\end{table}

In Table 4.12, the use case `MonitorLiveClassroom` renders student monitoring cards and logs on the web dashboard. The preconditions require that the exam session is ongoing. The trigger is the teacher opening live monitoring. The postconditions display live telemetry charts, and exception flows indicate offline states.

\begin{table}[H]
\centering
\caption{Use Case Table of Live Classroom Monitoring (MonitorLiveClassroom)}
\label{tab:uct_monitor_live_classroom}
\renewcommand{\arraystretch}{1.2}
\begin{tabular}{lp{10.6cm}}
\hline\hline
\textbf{Field} & \textbf{Details} \\
\hline\hline
\textbf{UCT ID} & UCT-12 \\
\textbf{Case Name} & MonitorLiveClassroom \\
\textbf{Overview} & Provides teachers with a live grid of student status cards, webcam feeds, and violation alerts. \\
\textbf{Actors} & Teacher, System Web Panel \\
\textbf{Pre-Conditions} & Students are connected and streaming telemetry. \\
\textbf{Triggers} & Teacher opens the Live Proctoring dashboard on the web console. \\
\textbf{Steps} & \begin{minipage}[t]{10.6cm}
1. Web panel connects to SignalR Hub and subscribes to exam channel.\\
2. Dashboard renders student monitoring cards in a grid layout.\\
3. System pushes live screenshot and focus logs to the matching cards.\\
4. Teacher views logs, alerts, and issues remote overrides if required.
\end{minipage} \\
\textbf{Post-Conditions} & Teacher dashboard is actively monitoring the lab terminals. \\
\textbf{Exception Flow} & \begin{minipage}[t]{10.6cm}
1. Server Disconnect: Dashboard displays warning and retries socket connection.\\
2. Inactive Student: Dashboard outlines student card in gray indicating offline state.
\end{minipage} \\
\hline\hline
\end{tabular}
\end{table}

In Table 4.13, the use case `SubmitExamAnswers` synchronizes student answers and removes client locks. The preconditions require that the student submits or the timer expires. The trigger is clicking submit. The postconditions write answers to the database and close the client workspace, with exception flows saving data locally.

\begin{table}[H]
\centering
\caption{Use Case Table of Controlled Exam Submission (SubmitExamAnswers)}
\label{tab:uct_submit_exam_answers}
\renewcommand{\arraystretch}{1.2}
\begin{tabular}{lp{10.6cm}}
\hline\hline
\textbf{Field} & \textbf{Details} \\
\hline\hline
\textbf{UCT ID} & UCT-13 \\
\textbf{Case Name} & SubmitExamAnswers \\
\textbf{Overview} & Coordinates final synchronization of answers, stops telemetry, and releases workstation locks. \\
\textbf{Actors} & Student, System \\
\textbf{Pre-Conditions} & Exam workspace is running and candidate answers are complete. \\
\textbf{Triggers} & Student clicks "Submit Exam" or exam duration timer expires. \\
\textbf{Steps} & \begin{minipage}[t]{10.6cm}
1. Client collects final answer state and queries server synchronization status.\\
2. Client transmits remaining answer packets to the API database.\\
3. System verifies completion status and terminates background logging threads.\\
4. Client releases OS lockdown hooks and closes the exam workspace.
\end{minipage} \\
\textbf{Post-Conditions} & Session is closed, student work is stored, and lockdown is released. \\
\textbf{Exception Flow} & \begin{minipage}[t]{10.6cm}
1. Network Disconnection: Client caches data locally and prompts student to consult invigilator.\\
2. Save Failures: System marks session pending and blocks workspace exit.
\end{minipage} \\
\hline\hline
\end{tabular}
\end{table}

In Table 4.14, the use case `EvaluateExamSubmissions` executes code in containers, semantic-grades theory, runs plagiarism checks, and exports PDF reports. The preconditions require that the exam session has concluded. The trigger is the teacher initiating evaluation. The postconditions save scores and plagiarism logs and download reports.

\begin{table}[H]
\centering
\caption{Use Case Table of Post-Exam Evaluation and Reporting (EvaluateExamSubmissions)}
\label{tab:uct_evaluate_exam_submissions}
\renewcommand{\arraystretch}{1.2}
\begin{tabular}{lp{10.6cm}}
\hline\hline
\textbf{Field} & \textbf{Details} \\
\hline\hline
\textbf{UCT ID} & UCT-14 \\
\textbf{Case Name} & EvaluateExamSubmissions \\
\textbf{Overview} & Runs compilations in Docker, semantic-grades theory, checks plagiarism, and exports reports. \\
\textbf{Actors} & System, Teacher, Administrator \\
\textbf{Pre-Conditions} & All student submissions are stored in the database. \\
\textbf{Triggers} & Teacher launches post-exam evaluation from the dashboard. \\
\textbf{Steps} & \begin{minipage}[t]{10.6cm}
1. System pulls code files and executes them in Docker containers against test cases.\\
2. System calls AI microservice to grade theory answers via cosine similarity.\\
3. System runs plagiarism checker to detect code copying across the class.\\
4. System compiles grades and calls QuestPDF to export downloadable reports.
\end{minipage} \\
\textbf{Post-Conditions} & Grades are updated, plagiarism flags are set, and PDF report is downloaded. \\
\textbf{Exception Flow} & \begin{minipage}[t]{10.6cm}
1. Compiler Sandbox Crash: System writes error traceback details to grade logs.\\
2. AI Service Timeout: System flags affected theory answers for manual grading.
\end{minipage} \\
\hline\hline
\end{tabular}
\end{table}

